% Problem record op_e432fc5ff73e63be. This TeX file is derived from
% database/problems_json/op_e432fc5ff73e63be.json by scripts/sync-tex.mjs; edit the JSON record
% and rerun the script rather than editing this file.
\section{Highest Clifford-hierarchy level of controlled gates with higher-level targets}
\label{sec:op_e432fc5ff73e63be}

\paragraph{Problem.}
What is the highest Clifford-hierarchy level that a controlled gate can
reach when its target lies in a given level $k\geq3$ and has a given Pauli
periodicity, and is that level bounded independently of the number of target
qubits?  Let $n\geq1$, and let $\mathcal P_n$ be the $n$-qubit Pauli group,
consisting of the operators $\omega P_1\otimes\cdots\otimes P_n$ with
$\omega\in\{\pm1,\pm i\}$ and $P_j\in\{I,X,Y,Z\}$.  The Clifford hierarchy is
defined by
\begin{equation}
  \mathcal C_1(n):=\mathcal P_n,
  \qquad
  \mathcal C_{k+1}(n):=\bigl\{U\in U(2^n):UPU^\dagger\in\mathcal C_k(n)
    \ \text{for every}\ P\in\mathcal P_n\bigr\},
  \qquad k\geq1.
  \label{eq:e432-hierarchy}
\end{equation}
Write $\mathcal C_\infty(n):=\bigcup_{k\geq1}\mathcal C_k(n)$, and let
$\ell_n(V):=\min\{k\geq1:V\in\mathcal C_k(n)\}$ be the exact level of
$V\in\mathcal C_\infty(n)$.  For $U\in U(2^n)$, define the controlled gate and
the Pauli periodicity
\begin{equation}
  CU:=|0\rangle\langle0|\otimes I+|1\rangle\langle1|\otimes U\in U(2^{n+1}),
  \qquad
  m(U):=\min\bigl\{t\in\mathbb Z_{\geq0}:U^{2^t}\in\mathcal P_n\bigr\},
  \label{eq:e432-periodicity}
\end{equation}
with $m(U):=\infty$ when no such $t$ exists.  Only the four phases
$\pm1,\pm i$ count as Pauli phases in Eqs.~\eqref{eq:e432-hierarchy}
and~\eqref{eq:e432-periodicity}.  This convention is essential: replacing $U$
by $e^{i\theta}U$ multiplies $CU$ by a phase gate on the control qubit, and it
can change both $m(U)$ and the level of $CU$.

For integers $k\geq2$ and $m\geq0$, the highest level reachable by controlling
a $k$th-level target with Pauli periodicity $m$ is
\begin{equation}
  L(n,k,m):=\max\bigl\{\ell_{n+1}(CU):
    U\in\mathcal C_k(n),\ m(U)=m,\ CU\in\mathcal C_\infty(n+1)\bigr\}.
  \label{eq:e432-max-level}
\end{equation}
The set in Eq.~\eqref{eq:e432-max-level} is nonempty because
$U=e^{i\pi/2^{m+1}}I$ qualifies, and the maximum is attained because
$\mathcal C_k(n)$ is finite modulo global phase while fixing $m(U)$ leaves
finitely many admissible phases.  Membership $U\in\mathcal C_k(n)$ includes the
lower levels, so $L(n,k,m)$ is nondecreasing in $k$.  Determine $L(n,k,m)$ for
all $n\geq1$, $k\geq3$, and $m\geq0$.  In particular, is $L(n,k,m)$ bounded by
a function of $k$ and $m$ alone, independently of the number $n$ of target
qubits?

\subsection*{Status}

Unsolved

\subsection*{Source}

Xu and Wang pose the question explicitly as Open Problem~1 in Section~5:
among all $n$-qubit unitaries in the $k$th level of the hierarchy with Pauli
periodicity $m$, what is the highest level that the controlled gate can reach?
They note that the answer can depend simultaneously on $m$, $n$, and $k$, and
that controlling non-Clifford targets could potentially bypass the exponential
qubit overhead inherent to Clifford targets
\sourcecite{ref:e432-xu-wang}{XW26}.  This record makes the discrete phase
convention of their Definition~4 and the requirement
$CU\in\mathcal C_\infty(n+1)$ explicit, and states the dependence on $n$ as a
boundedness question.

\subsection*{Progress}

\begin{itemize}
  \item Clifford targets are settled.  Xu and Wang proved the controlled-jump
  criterion (Theorem~3): every Clifford gate $U\notin\{\pm I\}$ with
  $m(U)<\infty$ satisfies
  \begin{equation}
    CU\in\mathcal C_{m(U)+2}(n+1)\setminus\mathcal C_{m(U)+1}(n+1)
    \label{eq:e432-clifford-jump}
  \end{equation}
  \sourcecite{ref:e432-xu-wang}{XW26}.  Their theorem is stated for all
  Clifford targets; the two excluded gates, for which $CU$ is a Pauli
  operator, matter only when $m(U)=0$.  For non-Pauli targets, Lemma~1 of
  Surti, Daguerre, and Kim had characterized the third level:
  $CU\in\mathcal C_3(n+1)$ exactly when $U$ is a Clifford gate with
  $U^2\in\mathcal P_n$ \sourcecite{ref:e432-surti-daguerre-kim}{SDK26}.
  Because $e^{i\pi/2^{m+1}}I$ is a Clifford gate with Pauli periodicity $m$,
  Eq.~\eqref{eq:e432-clifford-jump} gives $L(n,2,m)=m+2$ for all $n\geq1$ and
  $m\geq0$.

  \item The qubit count enters only after the global phase is discarded.
  Theorem~4 of Xu and Wang bounds the Pauli periodicity of an $n$-qubit
  Clifford gate by $\lceil\log_2(2n)\rceil$, and their Corollary~1 converts
  this into the requirement $n\geq2^{k-4}+1$ for a Clifford target whose
  controlled gate lies exactly in level $k\geq4$
  \sourcecite{ref:e432-xu-wang}{XW26}.  The proof bounds the order of the
  binary symplectic matrix of $U$, which ignores the global phase, whereas
  their Definition~4 of Pauli periodicity, like
  Eq.~\eqref{eq:e432-periodicity}, admits only the phases $\pm1,\pm i$.  Under
  that convention the bound fails as literally stated: $U=e^{i\pi/16}I$ is a
  one-qubit Clifford gate with $m(U)=3>\lceil\log_2 2\rceil$, and
  $CU=\operatorname{diag}(1,e^{i\pi/16})\otimes I$ lies exactly in level $5$
  on two qubits, in agreement with Eq.~\eqref{eq:e432-clifford-jump}.  The
  qubit-count bounds therefore hold for phase-normalized Clifford targets,
  whose global phase is chosen to minimize $m(U)$.  This caveat is a deduction
  made for this record, not a statement of the source.

  \item Every target obeys three lower bounds.  The repeated-commutator argument
  of Anderson and Weippert (Theorem~2.4 and Corollary~2.4.1), reproduced as the
  proof of Theorem~1 by Xu and Wang, uses no Clifford assumption:
  $CU\in\mathcal C_r(n+1)$ with $r\geq2$ forces $U^{2^{r-2}}\in\mathcal P_n$
  \sourcecite{ref:e432-anderson-weippert}{AW24},
  \sourcecite{ref:e432-xu-wang}{XW26}.  Separately, a direct calculation gives
  \begin{equation}
    CU\,(X\otimes I)\,CU^\dagger
    =(X\otimes I)\bigl(|0\rangle\langle0|\otimes U
      +|1\rangle\langle1|\otimes U^\dagger\bigr),
    \label{eq:e432-x-conjugation}
  \end{equation}
  and the block-diagonal lemma of Xu and Wang (Lemma~2) then places both $U$
  and $U^\dagger$ one level below $CU$.  Consequently, whenever
  $U\notin\{\pm I\}$ and $CU\in\mathcal C_\infty(n+1)$, the gates $U$ and
  $U^\dagger$ lie in the hierarchy and
  \begin{equation}
    \ell_{n+1}(CU)\geq\max\bigl\{m(U)+2,\ \ell_n(U)+1,\
      \ell_n(U^\dagger)+1\bigr\}.
    \label{eq:e432-lower-bound}
  \end{equation}
  The first term is the source argument; the second and third, obtained from
  Eq.~\eqref{eq:e432-x-conjugation}, are deductions made for this record.

  \item Non-Clifford targets already break the Clifford formula.  The
  classification of diagonal gates in the hierarchy (Theorem~3 of Cui,
  Gottesman, and Krishna), together with Eq.~\eqref{eq:e432-lower-bound},
  gives the exact levels, for $k\geq2$,
  \begin{equation}
    \begin{aligned}
      &U=C^{k-1}Z: && \ell_k(U)=k, && m(U)=1, && \ell_{k+1}(CU)=k+1;\\
      &U=\operatorname{diag}\bigl(1,e^{i\pi/2^{k-1}}\bigr): && \ell_1(U)=k,
        && m(U)=k-1, && \ell_2(CU)=k+1,
    \end{aligned}
    \label{eq:e432-diagonal-families}
  \end{equation}
  where $C^{k-1}Z$ is the $k$-qubit multi-controlled $Z$ gate
  \sourcecite{ref:e432-cui-gottesman-krishna}{CGK17}.  The first family in
  Eq.~\eqref{eq:e432-diagonal-families} exceeds $m(U)+2$ when $k\geq3$, and in
  both families the control raises the level by exactly one.  Anderson and
  Weippert (Section~3) ask whether this always happens, that is, whether
  $U\in\mathcal C_k\setminus\mathcal C_{k-1}$ implies
  $CU\in\mathcal C_{k+1}\setminus\mathcal C_k$
  \sourcecite{ref:e432-anderson-weippert}{AW24}.  Clifford targets already
  answer negatively: the three-qubit CNOT string whose linear part is a single
  Jordan block is a Clifford permutation with $m(U)=2$ by Corollary~2 of Xu and
  Wang, so $\ell_3(U)=2$ while $\ell_4(CU)=4$ by
  Eq.~\eqref{eq:e432-clifford-jump} \sourcecite{ref:e432-xu-wang}{XW26}.

  \item Third-level targets either leave the hierarchy under control or gain at
  least $k-1$ levels.  He, Robitaille, and Tan construct, for every $k\geq3$, a
  permutation gate $U_k$ on $2^k-1$ qubits with
  \begin{equation}
    U_k\in\mathcal C_3(2^k-1),
    \qquad
    U_k^{-1}\notin\mathcal C_k(2^k-1)
    \label{eq:e432-hrt-family}
  \end{equation}
  (Theorem~5.4 of the arXiv version, Theorem~42 of the proceedings version)
  \sourcecite{ref:e432-he-robitaille-tan}{HRT26}.  By
  Eqs.~\eqref{eq:e432-lower-bound} and~\eqref{eq:e432-hrt-family}, either
  $CU_k\notin\mathcal C_\infty(2^k)$ or $\ell_{2^k}(CU_k)\geq k+2$; neither
  alternative is established in the sources located.  If the second
  alternative holds for infinitely many $k$ along which $m(U_k)$ stays
  bounded, then $L(n,3,m)$ is unbounded in $n$ for some $m$, which answers the
  second question negatively.  The sources do not determine $m(U_k)$.  A
  direct computation for this record from the Toffoli-circuit form of $U_k$
  (Proposition~5.2 of the arXiv version) gives $m(U_3)=m(U_4)=2$, whereas
  $U_6^4$ is a nonidentity permutation fixing $|0\cdots0\rangle$, hence not a
  Pauli operator, so $m(U_6)\geq3$.  The dichotomy and these values are
  deductions made for this record.

  \item The condition $CU\in\mathcal C_\infty(n+1)$ in
  Eq.~\eqref{eq:e432-max-level} cannot be dropped.  Anderson and Weippert
  (Corollary~2.4.1) show that $CU\in\mathcal C_\infty(n+1)$ requires
  $U\in\mathcal C_\infty(n)$ and $m(U)<\infty$, and their Section~3 leaves
  sufficiency open \sourcecite{ref:e432-anderson-weippert}{AW24}.  De Silva
  and Lautsch give the explicit five-qubit gate
  $U=\operatorname{ctrl}_c(B)\operatorname{ctrl}_d(A)$, with control qubits
  $c,d$ and three-qubit Clifford targets $A=CX_{2\to3}CZ_{2,3}$ and
  $B=CX_{2\to3}CX_{1\to2}H_1$, and prove
  \begin{equation}
    U\in\mathcal C_5(5)\setminus\mathcal C_4(5),
    \qquad
    U^\dagger\notin\mathcal C_\infty(5)
    \label{eq:e432-dsl-gate}
  \end{equation}
  (Theorems~8.1 and~8.2 of their September 2026 preliminary preprint; their
  first level admits all global phases, which leaves every level $k\geq2$
  unchanged) \sourcecite{ref:e432-de-silva-lautsch}{dSL26}.  A direct matrix
  computation for this record gives $U^8=I$ and $U^4\notin\mathcal P_5$, so
  $m(U)=3$; Eqs.~\eqref{eq:e432-lower-bound} and~\eqref{eq:e432-dsl-gate} then
  give $CU\notin\mathcal C_\infty(6)$.  Thus hierarchy membership and finite
  Pauli periodicity of the target do not suffice.
\end{itemize}

\subsection*{References}

\begin{enumerate}[label={},leftmargin=4.8em,labelsep=0.5em]
  \footnotesize
  \item[\textup{[XW26]}]\label{ref:e432-xu-wang}
  Y. Xu and X. Wang, ``Controlled jump in the Clifford hierarchy,''
  arXiv:2602.22201 (2026).
  \href{https://doi.org/10.48550/arXiv.2602.22201}{doi:10.48550/arXiv.2602.22201};
  \href{https://arxiv.org/abs/2602.22201}{arXiv:2602.22201}.

  \item[\textup{[AW24]}]\label{ref:e432-anderson-weippert}
  J. T. Anderson and M. Weippert, ``Controlled gates in the Clifford
  hierarchy,'' arXiv:2410.04711 (2024), version 3 (2025).
  \href{https://doi.org/10.48550/arXiv.2410.04711}{doi:10.48550/arXiv.2410.04711};
  \href{https://arxiv.org/abs/2410.04711}{arXiv:2410.04711}.

  \item[\textup{[SDK26]}]\label{ref:e432-surti-daguerre-kim}
  S. Surti, L. Daguerre, and I. H. Kim, ``Efficient Simulation of Logical
  Magic State Preparation Protocols,'' \emph{PRX Quantum} \textbf{7}, 020329
  (2026).
  \href{https://doi.org/10.1103/fby6-xjbm}{doi:10.1103/fby6-xjbm};
  \href{https://arxiv.org/abs/2512.23799}{arXiv:2512.23799}.

  \item[\textup{[CGK17]}]\label{ref:e432-cui-gottesman-krishna}
  S. X. Cui, D. Gottesman, and A. Krishna, ``Diagonal gates in the Clifford
  hierarchy,'' \emph{Physical Review A} \textbf{95}, 012329 (2017).
  \href{https://doi.org/10.1103/PhysRevA.95.012329}{doi:10.1103/PhysRevA.95.012329};
  \href{https://arxiv.org/abs/1608.06596}{arXiv:1608.06596}.

  \item[\textup{[HRT26]}]\label{ref:e432-he-robitaille-tan}
  Z. He, L. Robitaille, and X. Tan, ``Characterization of Permutation Gates
  in the Third Level of the Clifford Hierarchy,'' in \emph{21st Conference on
  the Theory of Quantum Computation, Communication and Cryptography (TQC
  2026)}, LIPIcs \textbf{389}, 2:1--2:17 (2026).
  \href{https://doi.org/10.4230/LIPIcs.TQC.2026.2}{doi:10.4230/LIPIcs.TQC.2026.2};
  \href{https://arxiv.org/abs/2510.04993}{arXiv:2510.04993}.

  \item[\textup{[dSL26]}]\label{ref:e432-de-silva-lautsch}
  N. de Silva and O. Lautsch, ``The generalised semi-Clifford conjecture is
  false,'' arXiv:2609.11903 (2026), preliminary draft.
  \href{https://doi.org/10.48550/arXiv.2609.11903}{doi:10.48550/arXiv.2609.11903};
  \href{https://arxiv.org/abs/2609.11903}{arXiv:2609.11903}.
\end{enumerate}

\subsection*{Comment}

This is Open Problem~1 of Xu and Wang, with the phase convention and the
hierarchy-membership condition made explicit \sourcecite{ref:e432-xu-wang}{XW26}.
No determination of $L(n,k,m)$ for $k\geq3$ and no upper bound on the level of
a controlled gate with a non-Clifford target was located, and the works
located that cite Xu and Wang do not address it.  The Clifford rule in
Eq.~\eqref{eq:e432-clifford-jump}, the diagonal families in
Eq.~\eqref{eq:e432-diagonal-families}, and the lower bounds in
Eq.~\eqref{eq:e432-lower-bound} constrain $L(n,k,m)$ only from below.  A
complete answer requires matching upper bounds on $\ell_{n+1}(CU)$ for targets
in $\mathcal C_k(n)$ with $k\geq3$; a negative answer to the boundedness
question requires targets with fixed $k$ and $m$ whose controlled gates lie in
the hierarchy at levels unbounded in $n$.  The growth of the largest Pauli
periodicity attainable at a fixed level, Open Problem~2 of Xu and Wang, is the
separate record on \href{https://qiqc-op.com/problem/op_fe2c638b34652717/}{polynomial Pauli periodicity at a fixed
Clifford hierarchy level}.  Literature checked through 15 September 2026.

\subsection*{Field}

Quantum algorithm

\subsection*{Topic}

Clifford hierarchy

\subsection*{ID}

\texttt{op\_e432fc5ff73e63be}
